\section{Введение}

Задание лабораторной работы состоит в построении линейной регрессии для
прогнозирования непрерывного показателя успеваемости. Необходимо получить и
визуализировать описательную статистику набора данных, обработать пропуски и
категориальные признаки, нормировать числовые признаки, разделить наблюдения
на обучающую и тестовую части, самостоятельно найти коэффициенты методом
наименьших квадратов и оценить три модели с разными наборами признаков по
коэффициенту детерминации. Дополнительно исследуется синтетический признак.

Целевая переменная $y$ --- индекс успеваемости \textit{Performance Index}.
Исходные признаки обозначены так: $H$ --- часы учёбы, $P$ --- предыдущие
баллы, $E$ --- участие во внеучебных занятиях, $S$ --- часы сна, $Q$ ---
количество решённых пробных вариантов. Признак $E$ принимает значения
Yes или No; остальные признаки числовые. Индекс успеваемости округлён
до целого и лежит в пределах от 10 до 100.

Цель эксперимента --- установить, как меняется качество прогноза при
добавлении исходных и синтетического признаков. Данные являются
синтетическими, поэтому результаты характеризуют учебную задачу, а не
эффективность реального образовательного вмешательства.

\section{Описание метода}

\subsection{Линейная регрессия и метод наименьших квадратов}

Для $n$ наблюдений и $p$ подготовленных признаков множественная линейная
регрессия записывается в виде
\begin{equation}
\widehat y_i=\beta_0+\sum_{j=1}^{p}\beta_j z_{ij},
\qquad \widehat{\boldsymbol y}=X\boldsymbol\beta,
\end{equation}
где $z_{ij}$ --- значение $j$-го признака после преобразования, первый
столбец матрицы $X$ состоит из единиц, а $\beta_0$ --- свободный член.
Коэффициенты подбираются путём минимизации суммы квадратов остатков:
\begin{equation}
J(\boldsymbol\beta)=\sum_{i=1}^{n}(y_i-\widehat y_i)^2
=\|\boldsymbol y-X\boldsymbol\beta\|_2^2.
\end{equation}
Приравнивание производной к нулю приводит к системе нормальных уравнений
\begin{equation}
\nabla J=2X^{\mathsf T}(X\boldsymbol\beta-\boldsymbol y)=0,
\qquad
(X^{\mathsf T}X)\boldsymbol\beta=X^{\mathsf T}\boldsymbol y.
\end{equation}
В работе составляются $A=X^{\mathsf T}X$ и
$\boldsymbol b=X^{\mathsf T}\boldsymbol y$. Система
$A\boldsymbol\beta=\boldsymbol b$ решается методом Гаусса с выбором
максимального по модулю ведущего элемента в каждом столбце. После прямого
хода выполняется обратная подстановка. Обратная матрица явно не вычисляется.
Если ведущий элемент практически равен нулю, коэффициенты не определяются
однозначно и программа сообщает об ошибке.

При полном столбцовом ранге $X$ найденная точка является единственным
минимумом квадратичной функции $J$. Поскольку построение $X^{\mathsf T}X$
может ухудшать обусловленность, признаки предварительно стандартизируются;
для каждого решения дополнительно проверяется невязка нормальных уравнений.

\subsection{Подготовка данных}

Полностью одинаковые строки удаляются до разбиения. Это исключает попадание
одной и той же полной записи одновременно в обучение и тест. Категория
No кодируется числом $0$, категория Yes --- числом $1$. Для числовых пропусков
предусмотрено заполнение медианой обучающей части, для категориального
признака --- модой этой части. В исследуемом файле пропусков нет, однако
соответствующий этап остаётся частью алгоритма.

По обучающей части вычисляются среднее $\mu_j$ и стандартное отклонение
$s_j$. Затем к обеим частям применяется одно преобразование:
\begin{equation}
z_{ij}=\frac{x_{ij}-\mu_j^{\mathrm{train}}}
{s_j^{\mathrm{train}}},
\qquad
s_j^{\mathrm{train}}=
\sqrt{\frac{1}{n_{\mathrm{train}}}
\sum_{i\in\mathrm{train}}(x_{ij}-\mu_j^{\mathrm{train}})^2}.
\end{equation}
Для постоянного признака нулевой масштаб заменяется единицей. Параметры
заполнения и стандартизации не переоцениваются по тесту. Целевая переменная
не масштабируется.

\subsection{Оценка качества}

Основная метрика --- коэффициент детерминации на тестовой части:
\begin{equation}
R^2=1-\frac{\sum_i(y_i-\widehat y_i)^2}
{\sum_i(y_i-\overline y_{\mathrm{test}})^2}.
\end{equation}
Для характеристики размера ошибки также рассчитываются
\begin{equation}
\operatorname{RMSE}=\sqrt{\frac{1}{n}\sum_i(y_i-\widehat y_i)^2},
\qquad
\operatorname{MAE}=\frac{1}{n}\sum_i|y_i-\widehat y_i|.
\end{equation}
Больший $R^2$ и меньшие RMSE и MAE соответствуют лучшему прогнозу на
одинаковой выборке. Значение $R^2$ может быть отрицательным, если модель
хуже константного прогноза. Среднее тестовых значений в знаменателе
метрики не используется при обучении.

\section{Псевдокод метода}

\begin{enumerate}
\item Прочитать данные, удалить полные повторы и строки без целевого значения.
\item Перемешать строки с фиксированным зерном 42; отвести 80\% для обучения
и 20\% для теста.
\item По обучающей части определить значения заполнения, средние и масштабы
признаков; одинаково преобразовать обучение и тест.
\item Для каждого заранее заданного набора признаков дополнить матрицу $X$
столбцом единиц и сформировать $A=X^{\mathsf T}X$,
$\boldsymbol b=X^{\mathsf T}\boldsymbol y$.
\item Для каждого столбца $k$ выбрать в строках $k,\ldots,p$ элемент с
максимальным модулем; переставить строки и исключить переменную из нижних
строк. Обратной подстановкой найти $\boldsymbol\beta$.
\item Вычислить прогнозы для обучения и теста; рассчитать $R^2$, RMSE, MAE
и невязку $\|X^{\mathsf T}(X\boldsymbol\beta-\boldsymbol y)\|_\infty/n$.
\item Повторить расчёт для бонусного взаимодействия $H\cdot P$ и сопоставить
четыре варианта.
\end{enumerate}

Формирование нормальной системы требует порядка $O(np^2)$ операций,
её решение --- $O(p^3)$. Коэффициенты регрессии и метрики рассчитываются
с помощью арифметики NumPy; готовые регрессионные модели и библиотечные
решатели коэффициентов не применяются.

\section{Результаты выполнения}

\subsection{Состав и разделение данных}

Исходный набор содержит 10\,000 строк, пять признаков и целевой индекс.
Пропущенных значений не обнаружено. После удаления 127 полностью
совпадающих строк осталось 9\,873 наблюдения. Поскольку идентификаторов
студентов нет, одинаковые строки теоретически могут относиться к разным
объектам; их удаление является допущением эксперимента.

После перестановки с зерном 42 в обучающую часть попали 7\,898 строк,
в тестовую --- 1\,975. Для всех моделей используется одно и то же разбиение.
Заполнение пропусков, кодирование и стандартизация предусмотрены программой;
отсутствие пропусков означает, что фактические значения в этом эксперименте
не замещались.

\subsection{Описательная статистика и визуализация}

Основные числовые показатели приведены в таблице~\ref{tab:statistics}.
Стандартное отклонение в таблице рассчитано с делителем $n-1$, тогда как
масштаб стандартизации оценивается только на обучающей части с делителем
$n_{\mathrm{train}}$.
\begin{table}[H]\centering\small
\begin{tabular}{lrrrrr}\toprule
Показатель & $H$ & $P$ & $S$ & $Q$ & $y$\\\midrule
Количество & 9873 & 9873 & 9873 & 9873 & 9873\\
Среднее & 4.992 & 69.441 & 6.532 & 4.583 & 55.217\\
Станд. откл. & 2.589 & 17.326 & 1.698 & 2.867 & 19.209\\
Минимум & 1 & 40 & 4 & 0 & 10\\
Q1 (25\%) & 3 & 54 & 5 & 2 & 40\\
Медиана & 5 & 69 & 7 & 5 & 55\\
Q3 (75\%) & 7 & 85 & 8 & 7 & 70\\
Максимум & 9 & 99 & 9 & 9 & 100\\
\bottomrule\end{tabular}
\caption{Числовая статистика после удаления полных повторов.}
\label{tab:statistics}
\end{table}

\begin{figure}[H]
\centering\includegraphics[width=.98\textwidth]{distributions.png}
\caption{Распределения числовых и категориального признаков. Вертикальные
линии на гистограммах обозначают первый квартиль, медиану и третий квартиль.}
\label{fig:distributions}
\end{figure}
Предыдущие баллы лежат в интервале 40--99, индекс успеваемости --- 10--100.
По категориальному признаку наблюдаются 4\,986 значений No (50,50\%) и
4\,887 значений Yes (49,50\%). Число наблюдений, средние, отклонения,
экстремумы и квартили показаны также на рисунке~\ref{fig:distributions}.

\begin{figure}[H]
\centering\includegraphics[width=.68\textwidth]{correlation.png}
\caption{Парные корреляции после кодирования категориального признака.}
\label{fig:correlation}
\end{figure}
Наибольшая парная корреляция с целевой переменной наблюдается у предыдущих
баллов (около 0,92). Корреляция часов учёбы с индексом составляет около
0,38. Корреляционная матрица описывает данные, но не показывает причинное
влияние признаков.

\subsection{Модели и численные результаты}

В первой модели используется только $H$, во второй --- $H$ и $P$,
в третьей --- все пять исходных признаков $H,P,E,S,Q$. Для бонусной модели
к ним добавлено взаимодействие $HP=H\cdot P$, вычисленное до
стандартизации. Наборы признаков установлены заранее и не менялись
после просмотра тестовых оценок.

\begin{table}[H]\centering\small
\begin{tabular}{lrrrrr}\toprule
Модель & $p$ & $R^2$ train & $R^2$ test & RMSE test & MAE test\\\midrule
M1 & 1 & 0.140037 & 0.143777 & 17.9088 & 15.5060\\
M2 & 2 & 0.985786 & 0.985810 & 2.3055 & 1.8270\\
M3 & 5 & 0.988695 & 0.988620 & 2.0647 & 1.6397\\
M4 (бонус) & 6 & 0.988695 & 0.988619 & 2.0648 & 1.6398\\
\bottomrule\end{tabular}
\caption{Фактически полученные результаты на едином разбиении; $p$ не включает свободный член.}
\label{tab:metrics}
\end{table}


Константный прогноз, равный среднему индексу на обучении, даёт на тесте
$R^2=-0{,}000431$ и RMSE $=19{,}3583$. Все три основные модели превосходят
этот ориентир. Наиболее точная M3 достигает на тесте
$R^2=0{,}988620$, RMSE $=2{,}0647$ и MAE $=1{,}6397$.

\begin{figure}[H]
\centering\includegraphics[width=.95\textwidth]{comparison.png}
\caption{Тестовые значения $R^2$ и RMSE для трёх основных моделей и
бонусного варианта.}
\label{fig:comparison}
\end{figure}

После пересчёта коэффициентов M3 в исходные единицы получается
\begin{equation}
\widehat y=-34{,}12387+2{,}85139H+1{,}01888P+
0{,}62107E+0{,}48370S+0{,}19130Q.
\label{eq:m3}
\end{equation}
Например, коэффициент при $P$ показывает изменение условного прогноза
примерно на 1,019 при увеличении предыдущего результата на один балл и
неизменных остальных признаках. Значение свободного члена при нулевых
признаках находится вне наблюдаемой области и не имеет самостоятельного
педагогического смысла.

Корректность вычисления коэффициентов проверена на данных с известным
линейным решением, при перестановке ведущей строки и на вырожденной
матрице. Дополнительно проверены вычисление $R^2$ и заполнение
искусственно созданных пропусков без использования тестовой части.
Для каждой обученной модели максимальная нормированная невязка
$X^{\mathsf T}(X\boldsymbol\beta-\boldsymbol y)$ меньше
$9\cdot10^{-15}$.

\begin{figure}[H]
\centering\includegraphics[width=.95\textwidth]{diagnostics.png}
\caption{Фактические и прогнозные значения, а также остатки модели M3
на 1\,975 тестовых наблюдениях.}
\label{fig:diagnostics}
\end{figure}
Прогнозы сосредоточены около диагонали, а остатки колеблются вокруг нуля.
Такой график согласуется с небольшими средними ошибками, но сам по себе
не доказывает независимость или нормальность ошибок.

\section{Примеры использования метода}

Первые пять наблюдений в тестовой перестановке приведены в
таблице~\ref{tab:examples}. Они показывают, что ошибка одного прогноза
может отличаться от среднего значения метрики.
\begin{table}[H]\centering\small
\begin{tabular}{rrrrrrrrr}\toprule
Индекс & $H$ & $P$ & $E$ & $S$ & $Q$ & $y$ & $\widehat y$ & $y-\widehat y$\\\midrule
9483 & 3 & 90 & 1 & 8 & 2 & 74 & 71.003 & 2.997\\
2086 & 6 & 67 & 1 & 5 & 2 & 56 & 54.672 & 1.328\\
7700 & 7 & 61 & 1 & 4 & 7 & 55 & 51.883 & 3.117\\
9108 & 4 & 47 & 0 & 6 & 0 & 29 & 28.071 & 0.929\\
2773 & 4 & 49 & 1 & 6 & 7 & 32 & 32.069 & -0.069\\
\bottomrule\end{tabular}
\caption{Примеры M3 на тестовой части (исходные единицы).}
\label{tab:examples}
\end{table}


Для первого наблюдения $H=3$, $P=90$, $E=1$, $S=8$, $Q=2$.
Подстановка в уравнение~\eqref{eq:m3} даёт прогноз около 71,003 при
фактическом индексе 74; абсолютная ошибка составляет около 2,997.
Вторая строка показывает ошибку около 1,328. При использовании модели
для отдельного наблюдения необходимо сообщать, что RMSE описывает
совокупность тестовых ошибок, а не гарантирует точность каждого прогноза.

Линейная регрессия подходит для прозрачной базовой оценки непрерывного
показателя по нескольким измеренным факторам. Аналогичный подход может
использоваться при прогнозе расходов, спроса или иных количественных
величин, когда линейная зависимость является разумным приближением.
Интерпретируемые коэффициенты облегчают анализ, однако наблюдательная
связь не равнозначна причинному эффекту.

\section{Сравнение моделей}

При добавлении предыдущих баллов к часам учёбы тестовый $R^2$
увеличивается с 0,143777 у M1 до 0,985810 у M2. Следовательно, в данном
наборе именно $P$ даёт основное улучшение по сравнению с моделью только
по часам учёбы.

Использование всех пяти исходных признаков снижает RMSE с 2,3055 у M2
до 2,0647 у M3, то есть примерно на 10,45\%. Это показывает совместную
пользу $E$, $S$ и $Q$, но не позволяет приписать весь выигрыш одному
из этих признаков. Обучающий и тестовый $R^2$ M3 близки:
0,988695 и 0,988620 соответственно.

Добавление взаимодействия $HP$ почти не изменяет качество:
тестовый $R^2$ бонусной M4 равен 0,988619, RMSE --- 2,0648.
Обучающая ошибка немного уменьшается, а тестовая --- не улучшается.
По результату данного разбиения предпочтительна более простая M3.
Малое различие между M3 и M4 не следует трактовать как доказательство
статистического превосходства без повторных разбиений или внешней проверки.

\section{Заключение}

Реализована множественная линейная регрессия с самостоятельным решением
нормальных уравнений методом Гаусса. Выполнены описательная статистика,
визуализация, обработка повторов, кодирование, стандартизация, разделение
данных и сравнение трёх наборов признаков. Дополнительно проверено
взаимодействие часов учёбы с предыдущими баллами.

На фиксированной тестовой части лучшая основная модель M3, использующая
пять исходных признаков, получила $R^2=0{,}988620$,
RMSE $=2{,}0647$ и MAE $=1{,}6397$. Основной прирост качества связан
с добавлением предыдущих баллов; синтетический признак не дал заметного
улучшения. Результат относится к синтетическим данным и одному разбиению;
он не является оценкой причинных эффектов или гарантией точности на
реальных студенческих данных.
